\slajdid{09_3-s023}
\begin{frame}[label=09_3-s023]
Ono što sigurno treba da uočite jeste da će ovakav lanac kada nije opterećen parazitnom kapacitivnošću imati interno kašnjenje mnogo veće nego jedinični invertor. Ili ako mu se iz nekog razloga smanji parazitna kapacitivnost, neće imati optimalno kašnjenje.
\par\bigskip
Zbog toga su se u realizaciji CMOS logičkih kola koristila i \alert{\underline{nebaferisna CMOS logička kola}} (ako očekujemo da parazitne kapacitivnosti na izlazu nisu velike treba nam malo kašnjenje),
\par\bigskip
ali i \alert{\underline{duplo baferisana CMOS logička kola}} gde su na izlaz standardnog logičkog kola dodavana dva invertora čije su tranzistora povećane 4 odnosno 16 puta u odnosu na tranzistore u osnovnom logikom kolu, čime je predviđena neka standardna parazitna kapacitivnost na izlazu 64 puta veća od ulazne kapacitivnosti u logičko kolo.
\end{frame}
